\section{The local quaternion obstruction and the precise Euler input}
\label{sec:cubic-euler-prerequisites}

\begin{proposition}[AC-K29-split: ordinary inertia splits the Momose algebra]
\label{AC-K29-split}
Let $V$ be an actual type-II endpoint representation covered by
AC-K29-nonCM, and let $f$ be its modular form from AC-K29-typeII.
Let $F$ and $B/F$ be the inner-twist field and quaternion algebra
of $f$.  At the place $v\mid\ell$ corresponding to the chosen
coefficient embedding, $B_v=B\otimes_F F_v$ is split.
Consequently the characteristic-zero image hypotheses in
Lei--Loeffler--Zerbes, Proposition~7.1.4, hold for the restriction
of $V$ to $G_K$, for its dual, and for every scalar character twist
of either representation.  This assertion does not require a
residually irreducible lattice or integral big image.
\end{proposition}
\begin{proof}
We use the precise containment form of the classical image theorem:
Loeffler, \emph{Images of adelic Galois representations for modular
forms}, arXiv:1411.1789v2, \S2.2, Theorem~2.2.2, printed p.~6.
For the finite intersection $H$ of the kernels of the inner-twist
characters, it states that, in a suitable coefficient basis,
\[
                  \rho_{f,\ell}(H)\subset B_v^\times
                  \subset\mathrm{GL}_2(L_{\mathfrak P}).
                                                        \tag{CUE.1}
\]
Here $\rho_{f,\ell}$ is the cohomological normalization, with
geometric Frobenius traces; the source normalization of $f$ is
$V\otimes\kappa_\ell^{-e}$.  If the dual convention is chosen
instead, the same argument below takes inverses and transposes.

Intersect $H$ with $I_\ell$ and the kernel of the finite character
$\psi$.  This is an open subgroup $J$ of inertia.  By (G48.8),
the two diagonal characters of $V|_J$ are
$\kappa_\ell^e$ and $\kappa_\ell^{-e}$; the quadratic determinant
character is unramified at $\ell$.  The cyclotomic image of $J$
is open in $\mathbb Z_\ell^\times$.  Choose $g\in J$ with
$a=\kappa_\ell(g)$ and $a^{2e}\ne1$.
The two eigenvalues of $\rho_{f,\ell}(g)$ are then
$1$ and $a^{-2e}$, both in $\mathbb Q_\ell\subset F_v$, and
they are distinct.  This statement uses the triangular local
flag, so it does not assume that the local extension splits.

Put $x=\rho_{f,\ell}(g)\in B_v$.  Its reduced characteristic
polynomial is $(X-1)(X-a^{-2e})$.  If $B_v$ were a division
algebra, the equality
$(x-1)(x-a^{-2e})=0$ would force $x=1$ or $x=a^{-2e}$;
both scalar possibilities have repeated eigenvalues in the displayed
embedding, contradicting the two distinct eigenvalues just obtained.  A central simple algebra of degree
two over a local field is either division or a matrix algebra.
Therefore $B_v\simeq M_2(F_v)$.

For completeness, the remaining characteristic-zero group argument
is the proof of LLZ Proposition~7.1.4, printed pp.~25--26.
Their non-CM hypothesis is AC-K29-nonCM, and their Assumption~7.1.2
is now proved.  The image $\rho_{f,\ell}(G_K)$ of the
\emph{unnormalized} modular representation contains an open subgroup
of $\mathrm{GL}_2(\mathbb Z_\ell)$ after a coefficient basis change;
its derived image contains an open subgroup of
$\mathrm{SL}_2(\mathbb Z_\ell)$.  The Tate normalization agrees
with it on the derived group.  Thus the derived action on $V$ is absolutely
irreducible, it contains a nonidentity upper unipotent element
$\tau$ with $\dim_E V/(\tau-1)V=1$, and it contains
$\operatorname{diag}(y,y^{-1})$, $y\ne1$, with no fixed vector.
All scalar characters are trivial on the derived group.
Duality preserves these three assertions.  This is the
characteristic-zero conclusion, and gives no assertion that
the reduction of a chosen lattice contains
$\mathrm{SL}_2(\mathbb F_\ell)$.
\end{proof}

\paragraph{Exact source scope.}
Theorem~2.2.2 above is used as a classical theorem; its printed
proof is a reference to Momose and Ribet, with the ramification
bound supplied by Ghate--Gonzalez-Jimenez--Quer.  We have checked
the complete statement and that printed proof, and the complete
LLZ Proposition~7.1.4 proof.  We do not claim an independent
reproof of the upstream image theorem.  Its containment clause, rather than only the assertion that the image contains an open
subgroup, is essential to the argument.

\begin{proposition}[AC-K29-plus: an exact one-dimensional local comparison]
\label{AC-K29-plus}
Let $R$ be a two-dimensional de Rham representation over a
coefficient field $E$ of $G_K$, where $K$ is imaginary quadratic
and $\ell=\mathfrak p\overline{\mathfrak p}$ splits.
Put $U=R^*(1)$.  Suppose at each prime above $\ell$ that
\[
 H^0(K_v,U)=H^0(K_v,U^*(1))=0,\qquad
 \dim_E H_f^1(K_v,U)=1.                              \tag{CUE.2}
\]
Use Bloch--Kato local conditions away from $\ell$ throughout.
For $U$, let $\mathcal F$ impose $H_f^1$ at both primes, and
let $\mathcal F^+$ impose the full $H^1$ at $\mathfrak p$ and
$H_f^1$ at $\overline{\mathfrak p}$.  The dual structure
$(\mathcal F^+)^*$ on $R$ is zero at $\mathfrak p$ and
$H_f^1$ at $\overline{\mathfrak p}$.
If
\[
 H^1_{(\mathcal F^+)^*}(K,R)=0
 \quad\hbox{and}\quad
 z\in H^1_{\mathcal F^+}(K,U),\qquad
 \operatorname{loc}_{\mathfrak p}z\notin H_f^1(K_{\mathfrak p},U),
                                                        \tag{CUE.3}
\]
then $H_f^1(K,R)=0$.
\end{proposition}
\begin{proof}
Local Tate duality and the local Euler characteristic give
$\dim_EH^1(K_v,U)=2$.  Bloch--Kato local conditions on dual
representations are exact annihilators.  The comparison part of
Poitou--Tate duality for $\mathcal F\subset\mathcal F^+$ is
\[
 H^1_{\mathcal F^+}(K,U)\longrightarrow
 \frac{H^1(K_{\mathfrak p},U)}{H_f^1(K_{\mathfrak p},U)}
 \longrightarrow H_f^1(K,R)^*
 \longrightarrow H^1_{(\mathcal F^+)^*}(K,R)^* .        \tag{CUE.4}
\]
One may fix any common finite ramification set and impose the
same dual conditions at its other places; the displayed comparison
is unaffected.  The middle local quotient has dimension one.
The class $z$ makes the first arrow onto, and the last vector
space is zero.  Exactness forces $H_f^1(K,R)^*=0$.
\end{proof}

\paragraph{What the Euler-system source must additionally supply.}
For the actual target $R=V|_{G_K}\otimes\chi^j$, the symmetric
Hodge--Tate weights are $\{-e,e\}$.  With the LLZ convention that
$\kappa_\ell$ has weight $1$, those of $U=R^*(1)$ are
$\{1-e,1+e\}$ at \emph{each} prime above $\ell$.
Subject to the local invariant exclusions in (CUE.2), the
Bloch--Kato local dimension is consequently one at each prime.
The core-rank formula in LLZ Remark~B.2.4, printed pp.~40--41,
therefore gives core rank zero for $\mathcal F$ and one for
$\mathcal F^+$.  LLZ Theorem~7.4.2 uses a different infinity type,
with Bloch--Kato dimensions totaling three.  Its literal
everywhere-Bloch--Kato Euler class is not the correctly normalized
input for this symmetric target.

A sufficient coefficient-level Euler input is a lattice-bounded
rank-one Euler system for $U$ with Bloch--Kato local conditions
away from $\ell$, the ordinary positive-weight condition at
$\overline{\mathfrak p}$, and a nonzero singular projection at
$\mathfrak p$, together with the corresponding
characteristic-zero Euler-system bound for the dual strict group
in (CUE.3).  AC-K29-split supplies the image prerequisite for
that bound; (CUE.4) then supplies the complete local comparison.
This paragraph isolates an input: it does not assert that such
an Euler system has already been specialized to the actual
odd-weight, residually reducible forms.

In particular, the precise residual quotient condition in LLZ
Corollary~B.1.5 must be checked, or replaced by a proved uniform
bounded-error argument.  In the strict ordinary basis it concerns
the quotient character
$\kappa_\ell\delta^{-1}\chi^{-j}$ of $U$, not just $\chi^{-j}$.
The test $\chi(\mathrm{Frob}_\ell)\ne1$ alone does not remove
the other factors in that character.

\paragraph{Motivic interpretation and remaining specialization.}
The endpoint representation is the $\ell$-adic realization of
the modular-form motive in the source normalization, restricted
to $K$ and tensored with the cubic Artin character.  The dual
Tate twist in (CUE.4) is the realization-level duality pairing.
Rankin--Eisenstein classes are proposed motivic classes whose
\'etale realizations could furnish the indicated Euler system;
their existence with the necessary specialization, local
condition, uniform lattice bound and nonzero regulator is a
separate assertion.  B\"uy\"ukboduk--Lei, arXiv:1602.07508v3,
\S3.3.2 Proposition~3.10, printed pp.~19--21, constructs the
relevant kind of norm relations under the hypotheses retained
at the start of \S3.3.  The paper's global even-weight,
trivial-nebentype and residual-image assumptions are not silently
transferred to these odd-weight forms.  No Grothendieck archival
locator is claimed for this realization dictionary.
